\documentclass[12pt,a4paper]{ctexart}

\usepackage{amsmath,amssymb}
\usepackage{geometry}
\usepackage{enumitem}
\usepackage{hyperref}
\usepackage{booktabs}
\usepackage{listings}

\geometry{left=2.5cm,right=2.5cm,top=2.5cm,bottom=2.5cm}
\setlist{nosep,leftmargin=2em}

\title{第11章\quad 变化的磁场和变化的电场（电磁感应）}
\author{}
\date{}

\begin{document}

\maketitle

\section{知识点总结}

\subsection{电磁感应}

\begin{itemize}
    \item \textbf{电磁感应现象}：当穿过闭合导体回路的磁通量发生变化时，回路中出现感应电流。
    \item \textbf{磁通量}：$\Phi=\displaystyle\int_S\vec B\cdot\mathrm d\vec S=\int_S B\cos\theta\,\mathrm dS$；$\theta$ 为 $\vec B$ 与面积元 $\mathrm d\vec S$ 夹角。
    \item \textbf{三种使 $\Phi$ 变化的方式}：$B$ 变、$S$ 变、$\theta$ 变。
\end{itemize}

\subsection{感应电动势}

\begin{itemize}
    \item \textbf{法拉第电磁感应定律}：$\varepsilon=-\dfrac{\mathrm d\Phi}{\mathrm dt}$，负号体现楞次定律（感应电流方向反抗磁通变化）。
    \item \textbf{$N$ 匝密绕线圈}：$\varepsilon=-N\dfrac{\mathrm d\Phi}{\mathrm dt}=-\dfrac{\mathrm d(N\Phi)}{\mathrm dt}=-\dfrac{\mathrm d\Psi}{\mathrm dt}$，$\Psi=N\Phi$ 为磁通链。
    \item \textbf{感应电流}（回路电阻 $R$）：$I_i=\varepsilon/R=-\dfrac{1}{R}\dfrac{\mathrm d\Phi}{\mathrm dt}=\dfrac{\mathrm dq_i}{\mathrm dt}$。
    \item \textbf{感应电荷}：$q_i=\displaystyle\int_{t_1}^{t_2}I_i\,\mathrm dt=\dfrac{1}{R}(\Phi_1-\Phi_2)$（与过程无关）。
    \item \textbf{动生电动势}（磁场不变、导体运动）：非静电力是洛伦兹力。
    \begin{itemize}
        \item 电动势：$\varepsilon=\displaystyle\int_{-}^{+}(\vec v\times\vec B)\cdot\mathrm d\vec l$；$\mathrm d\varepsilon=Bv\,\mathrm dl\sin\theta$。
        \item 导体 $ab$ 在 $\vec B$ 中以 $\vec v$ 平动：$\varepsilon=\displaystyle\int_a^b(\vec v\times\vec B)\cdot\mathrm d\vec l$。
        \item \textbf{感应电动势做功来源}：外力克服安培力做功；洛伦兹力不做功（只传递能量）。
    \end{itemize}
    \item \textbf{感生电动势}（导体不动、磁场变化）：麦克斯韦提出\textbf{感生电场}（涡旋电场）$\vec E_V$，由变化磁场产生。
    \begin{itemize}
        \item $\varepsilon=\displaystyle\int_a^b\vec E_V\cdot\mathrm d\vec l$；闭合回路 $\varepsilon=\displaystyle\oint_L\vec E_V\cdot\mathrm d\vec l=-\dfrac{\mathrm d\Phi}{\mathrm dt}=-\displaystyle\int_S\dfrac{\partial\vec B}{\partial t}\cdot\mathrm d\vec S$。
        \item 微分形式：$\nabla\times\vec E_V=-\dfrac{\partial\vec B}{\partial t}$。
    \end{itemize}
    \item \textbf{感生电场性质}：无源有旋场（$\nabla\cdot\vec E_V=0$）；非保守场；与 $\dfrac{\partial\vec B}{\partial t}$ 成左手螺旋关系。
    \item \textbf{总感应电动势}：$\varepsilon=\displaystyle\int(\vec v\times\vec B)\cdot\mathrm d\vec l+\int\vec E_V\cdot\mathrm d\vec l$。
    \item \textbf{电子感应加速器条件}：轨道上 $B_R=\bar B/2$。
    \item \textbf{涡流}：变化磁场在导体内激起的涡旋状感应电流。
    \item \textbf{楞次定律}：感应电流的效果总是反抗引起感应电流的原因。
\end{itemize}

\subsection{自感和互感}

\begin{itemize}
    \item \textbf{自感现象}：线圈电流变化引起穿过自身磁通变化，从而产生自感电动势。
    \item \textbf{自感系数}：$\Psi=LI$，$\varepsilon_L=-\dfrac{\mathrm d\Psi}{\mathrm dt}=-L\dfrac{\mathrm dI}{\mathrm dt}-I\dfrac{\mathrm dL}{\mathrm dt}$，若 $L$ 恒定则 $\varepsilon_L=-L\dfrac{\mathrm dI}{\mathrm dt}$。
    \item \textbf{单位}：亨利 H = V$\cdot$s/A；自感 $L$ 由线圈匝数、几何、介质决定（线性介质中与 $I$ 无关）。
    \item \textbf{典型自感}：
    \begin{itemize}
        \item 长直螺线管：$L=\mu_0 N^2 S/l=\mu_0 n^2 V$。
        \item 同轴电缆：$L=\dfrac{\mu_0\mu_r l}{2\pi}\ln\dfrac{R_2}{R_1}$。
        \item 螺绕环（矩形截面）：$L=\dfrac{\mu_0 N^2 h}{2\pi}\ln\dfrac{R_2}{R_1}$。
    \end{itemize}
    \item \textbf{电磁惯性}：自感阻碍电流变化。
    \item \textbf{互感现象}：一线圈电流变化引起另一线圈磁通变化，产生互感电动势。
    \item \textbf{互感系数}：$\Psi_{21}=M_{21}I_1$，$\varepsilon_{21}=-M_{21}\dfrac{\mathrm dI_1}{\mathrm dt}$，$M_{12}=M_{21}=M$。
    \item \textbf{耦合系数}：$K=\dfrac{M}{\sqrt{L_1L_2}}\le 1$；$K=1$ 无漏磁。
    \item \textbf{线圈顺接}：$L=L_1+L_2+2M$；\textbf{反接}：$L=L_1+L_2-2M$。
    \item \textbf{互感典型}：
    \begin{itemize}
        \item 矩形线圈与长直导线（$N$ 匝）：$M=\dfrac{\mu_0 Nl}{2\pi}\ln\dfrac{b}{a}$。
        \item 同轴两圆线圈（$R\gg r,\ d\gg R$）：$M=\dfrac{\pi\mu_0 R^2 r^2}{2d^3}$。
    \end{itemize}
\end{itemize}

\subsection{磁场能量}

\begin{itemize}
    \item \textbf{自感磁能}：$W_m=\dfrac{1}{2}LI^2$。
    \item \textbf{磁场能量密度}：$w_m=\dfrac{B^2}{2\mu}=\dfrac{1}{2}\vec B\cdot\vec H$。
    \item \textbf{总磁场能}：$W_m=\displaystyle\int_V w_m\mathrm dV=\displaystyle\int_V\dfrac{1}{2}\vec B\cdot\vec H\,\mathrm dV$。
    \item \textbf{典型磁场能}：
    \begin{itemize}
        \item 长直螺线管：$W_m=\dfrac{B^2}{2\mu}V$。
        \item 螺绕环：$W_m=\dfrac{\mu_0\mu_r N^2 I^2 h}{4\pi}\ln\dfrac{R_2}{R_1}$。
        \item 同轴电缆（长 $l$）：$W_m=\dfrac{\mu I^2 l}{4\pi}\ln\dfrac{R_2}{R_1}$。
    \end{itemize}
    \item \textbf{与电场能对比}：$W_e=\dfrac{1}{2}CU^2$，$w_e=\dfrac{1}{2}\vec D\cdot\vec E$。
\end{itemize}

\subsection{麦克斯韦电磁场理论}

\begin{itemize}
    \item \textbf{位移电流}：$\vec j_D=\dfrac{\partial\vec D}{\partial t}$，$I_D=\dfrac{\mathrm d\Phi_D}{\mathrm dt}=\displaystyle\int_S\dfrac{\partial\vec D}{\partial t}\cdot\mathrm d\vec S$。
    \item \textbf{全电流}：$I_{\text{全}}=I_{\text{传导}}+I_{\text{位移}}$（全电流连续闭合）。
    \item \textbf{全电流安培环路定理}：$\displaystyle\oint_L\vec H\cdot\mathrm d\vec l=\displaystyle\int_S\left(\vec j+\dfrac{\partial\vec D}{\partial t}\right)\cdot\mathrm d\vec S$。
    \item \textbf{麦克斯韦方程组（积分形式）}：
    \begin{align}
        \oint_S\vec D\cdot\mathrm d\vec S &= \sum q_i=\int_V\rho\,\mathrm dV, \\
        \oint_S\vec B\cdot\mathrm d\vec S &= 0, \\
        \oint_L\vec E\cdot\mathrm d\vec l &= -\int_S\dfrac{\partial\vec B}{\partial t}\cdot\mathrm d\vec S, \\
        \oint_L\vec H\cdot\mathrm d\vec l &= \int_S\left(\vec j+\dfrac{\partial\vec D}{\partial t}\right)\cdot\mathrm d\vec S.
    \end{align}
    \item \textbf{麦克斯韦方程组（微分形式）}：
    \begin{align}
        \nabla\cdot\vec D &= \rho, & \nabla\cdot\vec B &= 0, \\
        \nabla\times\vec E &= -\dfrac{\partial\vec B}{\partial t}, & \nabla\times\vec H &= \vec j+\dfrac{\partial\vec D}{\partial t}.
    \end{align}
    \item \textbf{洛伦兹力公式}：$\vec F=q(\vec E+\vec v\times\vec B)$。
    \item \textbf{物质方程}：$\vec D=\varepsilon\vec E$，$\vec B=\mu\vec H$，$\vec j=\sigma\vec E$。
    \item \textbf{电磁场边值关系}：
    \begin{itemize}
        \item $\vec n\times(\vec E_2-\vec E_1)=0$；$\vec n\cdot(\vec D_2-\vec D_1)=\sigma$。
        \item $\vec n\times(\vec H_2-\vec H_1)=\vec j$；$\vec n\cdot(\vec B_2-\vec B_1)=0$。
    \end{itemize}
    \item \textbf{位移电流与传导电流比较}：
    \begin{itemize}
        \item 相同：都激发磁场。
        \item 不同：本质不同（位移电流为变化电场）；传导电流有热效应，位移电流无热效应；位移电流可在真空中。
    \end{itemize}
    \item \textbf{电磁波}（无源无旋波动方程）：$\nabla^2\vec E-\mu\varepsilon\dfrac{\partial^2\vec E}{\partial t^2}=0$。
    \begin{itemize}
        \item 介质中波速 $u=\dfrac{1}{\sqrt{\mu\varepsilon}}$，真空中 $c=\dfrac{1}{\sqrt{\mu_0\varepsilon_0}}=2.998\times10^8\,\text{m/s}$。
        \item 折射率 $n=\sqrt{\varepsilon_r\mu_r}=c/u$。
        \item 性质：$\vec E$、$\vec H$ 同相、同速；横波；$\sqrt{\varepsilon}E=\sqrt{\mu}H$；$\vec E\times\vec H$ 沿传播方向。
    \end{itemize}
    \item \textbf{能流密度（坡印廷矢量）}：$\vec S=\vec E\times\vec H$；电磁场能量密度 $w=\dfrac{1}{2}(\varepsilon E^2+\mu H^2)$。
\end{itemize}

\section{公式汇总}

\begin{enumerate}[leftmargin=3em,label=\textbf{公式\arabic*}：,ref=公式\arabic*]

    \item \textbf{磁通量}：$\Phi=\displaystyle\int_S\vec B\cdot\mathrm d\vec S$。

    \item \textbf{法拉第电磁感应定律}：$\varepsilon=-\dfrac{\mathrm d\Phi}{\mathrm dt}$；$N$ 匝 $\varepsilon=-N\dfrac{\mathrm d\Phi}{\mathrm dt}=-\dfrac{\mathrm d\Psi}{\mathrm dt}$。

    \item \textbf{感应电流与电荷}：$I_i=\varepsilon/R=-\dfrac{1}{R}\dfrac{\mathrm d\Phi}{\mathrm dt}$；$q_i=(\Phi_1-\Phi_2)/R$。

    \item \textbf{动生电动势}：$\varepsilon=\displaystyle\int(\vec v\times\vec B)\cdot\mathrm d\vec l$。
    \begin{itemize}
        \item 平直导线 $\varepsilon=Blv$；绕端转动铜棒 $\varepsilon=\dfrac{1}{2}B\omega R^2$；圆盘 $\varepsilon=\dfrac{1}{2}B\omega R^2$。
        \item 交流发电机 $\varepsilon=\varepsilon_m\sin\omega t$，$\varepsilon_m=NBS\omega$。
    \end{itemize}

    \item \textbf{感生电场与感生电动势}：$\varepsilon=\displaystyle\oint_L\vec E_V\cdot\mathrm d\vec l=-\int_S\dfrac{\partial\vec B}{\partial t}\cdot\mathrm d\vec S$，$\nabla\times\vec E_V=-\dfrac{\partial\vec B}{\partial t}$。
    \begin{itemize}
        \item 螺线管内：$E_V=\dfrac{r}{2}\dfrac{\partial B}{\partial t}$；外：$E_V=\dfrac{R^2}{2r}\dfrac{\partial B}{\partial t}$。
        \item 涡旋电场与 $\partial\vec B/\partial t$ 成左螺旋。
    \end{itemize}

    \item \textbf{自感}：$\Psi=LI$，$\varepsilon_L=-L\dfrac{\mathrm dI}{\mathrm dt}$。
    \begin{itemize}
        \item 长直螺线管：$L=\mu_0 n^2 V=\mu_0 N^2 S/l$。
        \item 螺绕环：$L=\dfrac{\mu_0 N^2 h}{2\pi}\ln\dfrac{R_2}{R_1}$。
        \item 同轴电缆：$L=\dfrac{\mu l}{2\pi}\ln\dfrac{R_2}{R_1}$。
    \end{itemize}

    \item \textbf{互感}：$\Psi_{21}=MI_1$，$\varepsilon_{21}=-M\dfrac{\mathrm dI_1}{\mathrm dt}$，$M_{12}=M_{21}=M$。
    \begin{itemize}
        \item 耦合系数：$K=\dfrac{M}{\sqrt{L_1L_2}}$。
        \item 顺接 $L=L_1+L_2+2M$；反接 $L=L_1+L_2-2M$。
    \end{itemize}

    \item \textbf{磁场能}：
    \begin{itemize}
        \item 自感磁能：$W_m=\dfrac{1}{2}LI^2$。
        \item 磁场能量密度：$w_m=\dfrac{1}{2}\vec B\cdot\vec H=\dfrac{B^2}{2\mu}$。
        \item 总磁场能：$W_m=\displaystyle\int_V\dfrac{1}{2}\vec B\cdot\vec H\,\mathrm dV$。
    \end{itemize}

    \item \textbf{位移电流}：$\vec j_D=\dfrac{\partial\vec D}{\partial t}$，$I_D=\dfrac{\mathrm d\Phi_D}{\mathrm dt}=\displaystyle\int_S\dfrac{\partial\vec D}{\partial t}\cdot\mathrm d\vec S$。

    \item \textbf{全电流安培环路定理}：$\displaystyle\oint_L\vec H\cdot\mathrm d\vec l=\displaystyle\int_S(\vec j+\partial\vec D/\partial t)\cdot\mathrm d\vec S$。

    \item \textbf{麦克斯韦方程组（积分）}：
    \begin{align}
        \oint_S\vec D\cdot\mathrm d\vec S &= \sum q_i, &
        \oint_S\vec B\cdot\mathrm d\vec S &= 0, \\
        \oint_L\vec E\cdot\mathrm d\vec l &= -\int_S\dfrac{\partial\vec B}{\partial t}\cdot\mathrm d\vec S, &
        \oint_L\vec H\cdot\mathrm d\vec l &= \int_S(\vec j+\partial\vec D/\partial t)\cdot\mathrm d\vec S.
    \end{align}

    \item \textbf{麦克斯韦方程组（微分）}：
    \[
        \nabla\cdot\vec D=\rho,\quad \nabla\cdot\vec B=0,\quad \nabla\times\vec E=-\dfrac{\partial\vec B}{\partial t},\quad \nabla\times\vec H=\vec j+\dfrac{\partial\vec D}{\partial t}.
    \]

    \item \textbf{物质方程}：$\vec D=\varepsilon\vec E$，$\vec B=\mu\vec H$，$\vec j=\sigma\vec E$。

    \item \textbf{边值关系}：$\vec n\times(\vec E_2-\vec E_1)=0$，$\vec n\cdot(\vec D_2-\vec D_1)=\sigma$，$\vec n\times(\vec H_2-\vec H_1)=\vec j$，$\vec n\cdot(\vec B_2-\vec B_1)=0$。

    \item \textbf{洛伦兹力}：$\vec F=q(\vec E+\vec v\times\vec B)$。

    \item \textbf{电磁波}：
    \begin{itemize}
        \item 波动方程：$\nabla^2\vec E-\mu\varepsilon\dfrac{\partial^2\vec E}{\partial t^2}=0$，$\nabla^2\vec H-\mu\varepsilon\dfrac{\partial^2\vec H}{\partial t^2}=0$。
        \item 介质中波速：$u=1/\sqrt{\mu\varepsilon}$；真空中 $c=1/\sqrt{\mu_0\varepsilon_0}=3\times10^8\,\text{m/s}$。
        \item 折射率：$n=\sqrt{\varepsilon_r\mu_r}=c/u$。
        \item 关系：$\sqrt{\varepsilon}E=\sqrt{\mu}H$；$\vec E$、$\vec H$、$\vec u$ 右手螺旋。
        \item 坡印廷矢量：$\vec S=\vec E\times\vec H$；电磁能密度 $w=\dfrac{1}{2}(\varepsilon E^2+\mu H^2)$。
    \end{itemize}

\end{enumerate}

\end{document}
